\medskip

L'espace est rapporté à un repère orthonormé \Oijk.

On considère les points 
\begin{center}A$(1~;~0~;~-1)$,\quad  B$(3~;~-1~;~2)$, \quad C$(2~;~-2~;~ -1)$ \quad et \quad D$(4~;~-1~;~-2)$.\end{center}

On note $\Delta$ la droite de représentation paramétrique 
\[\left\{\begin{array}{l c l}
x &=&\phantom{-}0\\
y &=&\phantom{-}2 + t\\
z &=& -1 + t
\end{array}\right.\:,\:\: \text{avec } t\in \R.\]

\begin{enumerate}
\item 
	\begin{enumerate}
		\item Montrer que les points A, B et C définissent un plan que l'on notera 
$\mathcal{P}$.
		\item Montrer que la droite (CD) est orthogonale au plan $\mathcal{P}$.
		
Sur le plan $\mathcal{P}$, que représente le point C par rapport à D ?
		\item Montrer qu'une équation cartésienne du plan $\mathcal{P}$ est: $2x + y - z - 3 = 0$.
	\end{enumerate}
\item 
	\begin{enumerate}
		\item Calculer la distance CD.
		\item Existe-t-il un point M du plan $\mathcal{P}$ différent de C vérifiant MD $= \sqrt 6$ ? Justifier la réponse.
	\end{enumerate}
\item 
	\begin{enumerate}
		\item Montrer que la droite $\Delta$ est incluse dans le plan $\mathcal{P}$.
		
Soit H le projeté orthogonal du point D sur la droite $\Delta$.
		\item Montrer que H est le point de $\Delta$ associé à la valeur $t = -2$ dans la représentation paramétrique de $\Delta$ donnée ci-dessus.
		\item En déduire la distance du point D à la droite $\Delta$.
	\end{enumerate}
\end{enumerate}

\bigskip
